\documentclass[11pt,letterpaper]{article}
\usepackage[T1]{fontenc}
\usepackage[utf8]{inputenc}
\usepackage[scaled=0.95]{helvet}
\renewcommand{\familydefault}{\sfdefault}
\usepackage[top=1.7cm,bottom=1.5cm,left=2.2cm,right=2.2cm]{geometry}
\usepackage{microtype}
\usepackage[hidelinks]{hyperref}
\pagestyle{empty}
\setlength{\parindent}{0pt}
\setlength{\parskip}{7pt}

\begin{document}

{\large\bfseries Grover Aruquipa}\par\vspace{-6pt}
{\small Qu\'ebec, QC, Canada \ \textbullet\ \ grover.aruquipa.9@gmail.com \ \textbullet\ \ groveraruquipa.github.io}

\vspace{4pt}
Hiring Team, MacLean Engineering, Ontario\\
Team Lead, Vehicle Automation (Applications)

Dear Hiring Team,

An underground mine is one of the hardest places to make a machine work on its own: no GPS, dust, poor communications, heavy equipment and people nearby. Taking an automation feature from a prototype to something a mine trusts every shift is the kind of problem I want to own, and MacLean is one of the few companies in Canada building that autonomy on its own vehicles.

I am finishing my PhD at Universit\'e Laval with Prof.\ Cl\'ement Gosselin and Prof.\ Patrice Lambert, and I defend in early 2027. Over three years I led three robotic systems with 4, 7 and 8 actuators from concept to validated operation. On each one I wrote the full software chain myself: the kinematic models, the planning with collision and singularity checks, and the real-time control in C++ and Simulink at 500\,Hz on OPAL-RT. I then integrated it with the motors, drives and encoders and tested it on the hardware until it met its targets. I build autonomy software in ROS~2 and MoveIt~2, and I simulate in Isaac Sim and MuJoCo before anything moves for real, which is the same simulation-first loop your team uses to reach the vehicle.

I lead people as well as code. I am the technical lead of two graduate engineers on an 8-actuator robot, from concept to prototype, and I led the Laval team that won first place at the ASME Student Mechanisms and Robotics Design Competition in 2025, as well as our team at a MassRobotics final in Boston. I split the work into parts others can own, I review it, and I bring it back together on the bench. Before the PhD I worked in industry as a computer vision engineer at Sencon, building vision systems for industrial sites, and as a data science intern at Hitachi ABB Power Grids in Zurich, applying machine learning to sensor data from high-voltage equipment. Perception and telemetry from real equipment are not new to me.

I also speak Spanish natively and French fluently enough to work in it every day in Qu\'ebec. With MacLean's customers and operations in Latin America and Mexico, I can work directly with mine sites and colleagues there in their own language.

I would be glad to lead the Vehicle Automation Applications team, and I would be equally glad to join it as a senior automation engineer if that is where your plans place me. I am available from January 2027 and ready to relocate within Ontario and work in the office. Prof.\ Gosselin, Prof.\ Lambert and Univ.-Prof.\ Andreas M\"uller of JKU Linz can speak to my work; their details are in my CV.

Thank you for your time. I would be happy to walk you through my robots and my code.

Sincerely,\\
Grover Aruquipa

\end{document}
